% =====================================================================
%  PREÁMBULO — paquetes y configuración tipográfica
%  Normalmente no hace falta tocar este fichero. Si añades un paquete,
%  hazlo ANTES del bloque final (hyperref/cleveref deben ir los últimos).
% =====================================================================

% Evita que pdfTeX incruste rutas locales de los ficheros incluidos en el PDF
\ifdefined\pdfsuppressptexinfo \pdfsuppressptexinfo=-1 \fi

% --- Codificación, idioma y tipografía ---------------------------------
\usepackage[utf8]{inputenc}
\usepackage[T1]{fontenc}
\usepackage{lmodern}
\usepackage[english, spanish, es-tabla, es-noshorthands]{babel}  % el último idioma es el principal
% Números romanos de la parte preliminar en minúscula normal (babel-spanish
% los pone en versalitas, que no existen en negrita y generan avisos)
\makeatletter\def\es@scroman#1{#1}\makeatother
\usepackage{microtype}
\usepackage{csquotes}
\usepackage{etoolbox}

% --- Página -----------------------------------------------------------
\usepackage[a4paper, top=2.5cm, bottom=2.5cm, inner=3cm, outer=2.5cm,
            headheight=15pt]{geometry}
\usepackage{emptypage}          % páginas de relleno realmente en blanco
\usepackage{fancyhdr}
\raggedbottom

% Desaconseja (sin prohibir) partir palabras al final de línea
\hyphenpenalty=5000
\exhyphenpenalty=5000

% --- Matemáticas y unidades -------------------------------------------
\usepackage{amsmath, amssymb, mathtools}
\usepackage{siunitx}
\sisetup{output-decimal-marker={,}, per-mode=symbol}
\usepackage{eurosym}

% --- Figuras, tablas y dibujos ----------------------------------------
\usepackage{graphicx}
\graphicspath{{figuras/}}
\usepackage{float}
\usepackage[labelfont=bf]{caption}
\usepackage{subcaption}
\usepackage{booktabs, tabularx, multirow}
\usepackage[table]{xcolor}
\usepackage{lscape}
\usepackage{pdfpages}
\usepackage{tikz}
\usetikzlibrary{arrows.meta, positioning, fit, calc, shapes.geometric, backgrounds}
\usepackage{pgfplots}
\pgfplotsset{compat=1.18}
\usepackage{pgfgantt}

% --- Código y algoritmos -----------------------------------------------
\usepackage{listings}
\definecolor{codegreen}{rgb}{0,0.6,0}
\definecolor{codegray}{rgb}{0.5,0.5,0.5}
\definecolor{codepurple}{rgb}{0.58,0,0.82}
\definecolor{backcolour}{rgb}{0.95,0.95,0.92}
\lstdefinestyle{memoria}{
    backgroundcolor=\color{backcolour},
    commentstyle=\color{codegreen},
    keywordstyle=\color{codepurple},
    numberstyle=\tiny\color{codegray},
    stringstyle=\color{codepurple},
    basicstyle=\ttfamily\small,
    breaklines=true,
    captionpos=b,
    keepspaces=true,
    numbers=left,
    showstringspaces=false,
    tabsize=4,
    frame=single,
    literate={á}{{\'a}}1 {é}{{\'e}}1 {í}{{\'i}}1 {ó}{{\'o}}1 {ú}{{\'u}}1
             {ñ}{{\~n}}1 {Á}{{\'A}}1 {É}{{\'E}}1 {Í}{{\'I}}1 {Ó}{{\'O}}1
             {Ú}{{\'U}}1 {Ñ}{{\~N}}1 {ü}{{\"u}}1
}
\lstset{style=memoria}
\usepackage[ruled,vlined,linesnumbered,algochapter,spanish,onelanguage]{algorithm2e}
\SetAlgorithmName{Algoritmo}{algoritmo}{Índice de algoritmos}

% --- Anexos, acrónimos y bibliografía ------------------------------------
\usepackage[titletoc,title]{appendix}
\usepackage[printonlyused]{acronym}  % solo lista las siglas usadas
\usepackage[backend=biber, style=numeric, sorting=none]{biblatex}
\addbibresource{bibliografia/referencias.bib}

% --- Nombres en español ------------------------------------------------
\addto\captionsspanish{%
  \renewcommand{\appendixname}{Anexo}%
  \renewcommand{\lstlistingname}{Código}%
  \renewcommand{\lstlistlistingname}{Índice de código}%
  \renewcommand{\bibname}{Bibliografía}%
}
\renewcommand{\appendixtocname}{Anexos}
\renewcommand{\appendixpagename}{Anexos}

% --- Cabeceras y pies ----------------------------------------------------
\fancyhf{}
\fancyhead[LE]{\nouppercase{\slshape\leftmark}}
\fancyhead[RO]{\nouppercase{\slshape\rightmark}}
\fancyfoot[LE,RO]{\thepage}
\renewcommand{\headrulewidth}{0.4pt}
% Consejo: si un título no cabe en la cabecera, usa \chapter[corto]{largo}

% --- Enlaces y referencias cruzadas (SIEMPRE AL FINAL) --------------------
\usepackage[hidelinks]{hyperref}
\hypersetup{
  pdftitle={\Titulo},
  pdfauthor={\Autor},
  pdfsubject={\TipoTrabajo. \Titulacion},
  pdfkeywords={\PalabrasClave},
}
\usepackage[spanish, noabbrev]{cleveref}
\crefname{table}{tabla}{tablas}
\Crefname{table}{Tabla}{Tablas}
\crefname{appendix}{anexo}{anexos}
\Crefname{appendix}{Anexo}{Anexos}
\crefname{lstlisting}{código}{códigos}
\Crefname{lstlisting}{Código}{Códigos}
\crefname{algocf}{algoritmo}{algoritmos}
\Crefname{algocf}{Algoritmo}{Algoritmos}
